În timpul procesului de imprimare 3D pot apărea diverse defecte care pot cauza pierderi de material, pierderi de timp și chiar deteriorarea echipamentului.
Scopul acestei lucrări este dezvoltarea unui sistem de monitorizare în timp real, capabil să recunoască defectele de imprimare folosind metode de procesare a imaginilor.
Soluția propusă se bazează pe un sistem de supraveghere cu mai multe camere, care înregistrează continuu procesul de imprimare, iar cadrele sunt
prelucrate de o rețea neuronală YOLOv8 care rulează într-un mediu bazat pe cloud. Modelul a fost antrenat pe un set de date adnotat propriu și este capabil
să recunoască 9 clase de defecte, atingând o precizie de 89.2\% mAP50 și o valoare recall de 83.6\%, iar prelucrarea fiecărui cadru necesită aproximativ 23 ms timp de inferență.
Sistemul include și un mecanism de notificare, care oferă utilizatorului un feedback imediat în cazul detectării unei anomalii.
Pe baza rezultatelor experimentale, sistemul este potrivit pentru detectarea defectelor în timp real și poate contribui la creșterea eficienței procesului de
imprimare, precum și la reducerea pierderilor de material.
